\section{Results}

This section evaluates how generated tabular lens rows become usable optical prescriptions. Unless otherwise stated, each result starts from 10,000 generated samples. We report the \emph{simulation-verified feasible sample fraction}: the percentage of generated samples that pass post-processing and CodeV/LLESP preparation and then return from CodeV without an error. This denominator is stricter than a success rate on submitted samples because it also counts samples rejected before simulation. Log-whitened material-parameterization results and CoDi runs are excluded from the main results in this version.

\subsection{Synthetic-Expanded Training Data Improves Feasible Generation}
\label{sec:results_synthetic_expansion}

The first experiment asks whether the synthetic-expanded training set improves usable generation under the final pipeline. To isolate this effect, the representation is fixed to the tuned parameterization, NRTE post-processing is enabled, and each model is compared between patent-only training and the patent-plus-synthetic 250k training variant. Model hyperparameters are held fixed within each model family; the intended variable in this subsection is the training data source, not model size, representation, or post-processing.

\begin{center}
\begin{minipage}{0.98\linewidth}
\centering
\footnotesize
\captionof{table}{\textbf{Patent-only versus patent-plus-synthetic 250k feasible generation.} Deltas are absolute percentage-point changes from patent-only to 250k training. The dominant remaining failure mode is computed on the 250k run using mutually exclusive failure categories.}
\label{tab:section_6_1_synthetic_expansion_summary}
\begin{tabular}{lrrrl}
\toprule
\textbf{Model} & \multicolumn{2}{c}{\textbf{Feasible fraction $\uparrow$ (\%)}} & \textbf{$\Delta$ feasible} & \textbf{Dominant 250k failure} \\
\cmidrule(lr){2-3}
 & \textbf{Patent only} & \textbf{Patent+synth 250k} & \textbf{(pp)} & \\
\midrule
SMOTE-style reference & 81.40 & 98.43 & +17.03 & CodeV error \\
STaSy & 1.70 & 58.53 & +56.83 & Material only \\
TabDDPM & 0.00 & 0.00 & +0.00 & Material + geometry \\
TabSyn & 0.57 & 77.48 & +76.91 & Material only \\
\bottomrule
\end{tabular}
\end{minipage}
\end{center}

Table~\ref{tab:section_6_1_synthetic_expansion_summary} shows the feasible fraction for the four matched model families. The synthetic-expanded 250k dataset substantially increases the end-to-end feasible fraction for the learned tabular generators that can use the added data. TabSyn increases from 0.57\% feasible samples under patent-only training to 77.48\% under patent-plus-synthetic 250k training, and STaSy increases from 1.70\% to 58.53\%. The SMOTE-style reference is already highly feasible on patent-only training, but still increases from 81.40\% to 98.43\%. TabDDPM remains at 0.00\% feasible samples in both settings, showing that the benefit of synthetic-expanded training is not automatic across all tabular generators.

These results support two conclusions. First, the synthetic-expanded dataset is useful because the patent-only set is too small for several learned generators to produce simulation-ready prescriptions at meaningful rates. Second, feasibility must be measured after the optical reconstruction and simulation-readiness pipeline. A generator can produce rows that look tabularly valid while still failing material reconstruction, geometric consistency checks, or CodeV execution. The mutually exclusive failure-mode breakdown in Appendix~\ref{app:section_6_1_failure_breakdown} shows that the 250k runs reduce the large material-plus-geometry failure mass seen in the patent-only learned-model runs, but the remaining failure modes differ by generator: material-only failures dominate for STaSy and TabSyn, CodeV errors dominate the small residual failure rate for SMOTE, and TabDDPM remains dominated by coupled material and geometry failures.

\subsection{Effect Of Tuned Optical Representation}

The second experiment isolates the effect of the model-facing representation. The dataset is fixed to patent-plus-synthetic 250k, NRTE post-processing is enabled, and the comparison is restricted to model families with matched tuned and untuned returned runs. This keeps the intended variable as the optical representation rather than dataset size or post-processing.

Table~\ref{tab:tuned_untuned_results} shows that the tuned representation has little effect on the SMOTE-style reference but a large effect on TabSyn. For SMOTE, both representations submit nearly all generated rows to CodeV/LLESP and return nearly all submitted rows successfully, giving only a 0.36 percentage-point feasible-yield change. For TabSyn, the tuned and untuned branches submit nearly the same fraction of rows, 85.77\% versus 86.06\%, but their CodeV success rates among submitted rows differ sharply: 90.33\% for tuned versus 59.12\% for untuned. The resulting simulation-verified feasible fraction increases from 50.88\% to 77.48\%, a 26.60 percentage-point gain. This indicates that the tuned representation primarily improves the quality of reconstructed optical prescriptions that reach simulation, not merely the number of rows that pass the pre-CodeV gate.

\begin{center}
\begin{minipage}{0.98\linewidth}
\centering
\scriptsize
\setlength{\tabcolsep}{3pt}
\captionof{table}{\textbf{Tuned versus untuned representation comparison on patent-plus-synthetic 250k runs.} All entries use NRTE post-processing and the same 10,000-sample denominator. Deltas are tuned minus untuned in percentage points.}
\label{tab:tuned_untuned_results}
\begin{tabular}{lrrrrrrr}
\toprule
\textbf{Model} & \multicolumn{2}{c}{\textbf{CodeV input accepted (\%)}} & \multicolumn{2}{c}{\textbf{Success on submitted (\%)}} & \multicolumn{2}{c}{\textbf{Feasible fraction $\uparrow$ (\%)}} & \textbf{$\Delta$ feasible} \\
\cmidrule(lr){2-3}\cmidrule(lr){4-5}\cmidrule(lr){6-7}
 & \textbf{Tuned} & \textbf{Untuned} & \textbf{Tuned} & \textbf{Untuned} & \textbf{Tuned} & \textbf{Untuned} & \textbf{(pp)} \\
\midrule
SMOTE-style reference & 99.76 & 99.82 & 98.67 & 98.25 & 98.43 & 98.07 & +0.36 \\
TabSyn & 85.77 & 86.06 & 90.33 & 59.12 & 77.48 & 50.88 & +26.60 \\
\bottomrule
\end{tabular}
\end{minipage}
\end{center}

\subsection{Effect Of Synthetic Dataset Scale}

To isolate dataset scale from the representation ablation, Table~\ref{tab:dataset_scale_results} fixes the generator to TabSyn, the parameterization to the tuned representation, and the post-processing branch to NRTE. Since the paper-level question is how dataset scale changes the final usable sample population, the table reports the end-to-end feasible yield and a mutually exclusive breakdown of infeasible rows rather than intermediate pipeline gates.

The patent-only row is retained only as the low-data anchor. The main dataset effect is the transition from patent-only to patent-plus-synthetic training, where the feasible yield rises from 0.57\% to 77.48\%. Beyond 250k synthetic-expanded training rows, however, this fixed TabSyn configuration does not improve monotonically: the 500k run is comparable to 250k, while the full-dataset run falls to 67.15\%. The full-dataset run has higher material-only, material-plus-geometry, and CodeV-error fractions than the 250k run, so the drop is visible directly in the final infeasibility composition. Thus, the million-scale dataset is best interpreted as a released scaling endpoint and resource for higher-capacity or retuned models, not as a guarantee that the same off-the-shelf TabSyn configuration will be optimal at the largest scale.

\begin{center}
\begin{minipage}{0.98\linewidth}
\centering
\scriptsize
\setlength{\tabcolsep}{4.5pt}
\captionof{table}{Dataset-scale comparison for TabSyn with tuned representation and NRTE post-processing. Percentages are computed over 10{,}000 generated rows. Infeasibility categories are mutually exclusive and sum to the any-infeasible column.}
\label{tab:dataset_scale_results}
\begin{tabular}{lrrrrrr}
\toprule
\textbf{Training data} & \shortstack{\textbf{Feasible}\\\textbf{yield $\uparrow$}} & \shortstack{\textbf{Any}\\\textbf{infeasible $\downarrow$}} & \shortstack{\textbf{Material}\\\textbf{only $\downarrow$}} & \shortstack{\textbf{Geometry}\\\textbf{only $\downarrow$}} & \shortstack{\textbf{Material}\\\textbf{+ geom. $\downarrow$}} & \shortstack{\textbf{CodeV}\\\textbf{error $\downarrow$}} \\
\midrule
Patent only & 0.57 & 99.43 & 35.00 & 0.01 & 59.72 & 4.70 \\
Patent+synth 250k & 77.48 & 22.52 & 11.75 & 0.00 & 2.48 & 8.29 \\
Patent+synth 500k & 75.97 & 24.03 & 11.22 & 0.00 & 2.76 & 10.05 \\
Patent+synth full & 67.15 & 32.85 & 15.75 & 0.00 & 4.59 & 12.51 \\
\bottomrule
\end{tabular}
\end{minipage}
\end{center}

\subsection{Ray-Tracing Based Post-Processing Benefit}
\label{sec:results_raytrace_postproc}

The final ablation evaluates the optional ray-tracing based post-processing layer. The generator, representation, and training data are fixed to TabSyn with the tuned representation on the patent-plus-synthetic 250k training variant. The only intended variable is whether generated prescriptions pass directly to CodeV/LLESP preparation or first pass through NRTE.

NRTE increases the end-to-end feasible yield from 66.52\% to 77.48\%, a gain of 10.96 percentage points under the same 10{,}000-sample denominator (Table~\ref{tab:raytrace_postproc_benefit}). This shows that the deterministic correction layer adds simulation-verified usable designs without retraining the generator.

\begin{center}
\begin{minipage}{0.98\linewidth}
\centering
\footnotesize
\setlength{\tabcolsep}{6pt}
\captionof{table}{\textbf{Effect of ray-tracing based post-processing on TabSyn tuned 250k generated samples.} Delta is NRTE minus no-NRTE in percentage points, with the same 10{,}000-sample denominator used for both branches.}
\label{tab:raytrace_postproc_benefit}
\begin{tabular}{lrrr}
\toprule
\textbf{Dataset} & \multicolumn{2}{c}{\textbf{Feasible fraction $\uparrow$ (\%)}} & \textbf{$\Delta$ feasible $\uparrow$ (pp)} \\
\cmidrule(lr){2-3}
 & \textbf{No NRTE} & \textbf{NRTE} & \textbf{NRTE--no NRTE} \\
\midrule
Patent+synth 250k & 66.52 & 77.48 & +10.96 \\
\bottomrule
\end{tabular}
\end{minipage}
\end{center}

Feasibility alone does not show whether the correction damages optical quality. Figure~\ref{fig:results_nrte_metric_improvement_distribution_250k} therefore compares paired optical metrics on common rows that were submitted in both branches and returned successfully in both CodeV runs. The plotted value is $\log_{10}(\mathrm{no\ NRTE}/\mathrm{NRTE})$ for each field-metric observation; positive values indicate lower metric values after NRTE. Spot-size and encircled-energy metrics are normalized by the corresponding Airy diameter before comparison, while the nominal performance metric is left in its returned units. This paired comparison is a within-sample post-processing diagnostic rather than a ranking of unrelated generated designs.

\begin{center}
    \includegraphics[width=\linewidth]{../../../code/optics_data_modder_2/results/figures/paper_section_6_4_paired_metric_improvements_250k_20260522/section_6_4_paired_metric_improvement_violin_250k.pdf}
    \captionof{figure}{\textbf{Paired optical-metric improvement distributions after NRTE for TabSyn tuned 250k generations.} Each distribution pools the five evaluated field positions over common rows that succeeded in both CodeV/LLESP branches. Positive values indicate lower metric values after NRTE. SPO67, SPO100, EE50, and EE90 are normalized by the Airy diameter; nominal performance is not Airy-normalized.}
    \label{fig:results_nrte_metric_improvement_distribution_250k}
\end{center}

Table~\ref{tab:raytrace_postproc_quality_250k} reports the corresponding median values and paired improvement fractions. NRTE reduces the median normalized SPO67 from 168.24 to 22.95, normalized SPO100 from 285.18 to 57.59, and nominal performance from 11.76 to 2.73. The median reductions for EE50 and EE90 are smaller but remain positive, and 73.2--75.9\% of valid field-metric observations improve. Field-resolved median reductions are provided in Appendix~\ref{app:raytrace_metric_delta_details}.

\begin{center}
\begin{minipage}{0.98\linewidth}
\centering
\scriptsize
\setlength{\tabcolsep}{4pt}
\captionof{table}{\textbf{Paired optical-quality summary for TabSyn tuned 250k samples with and without NRTE.} $N$ is the number of valid paired field-metric observations after requiring common rows, successful CodeV/LLESP returns in both branches, and positive finite metric values. Percent reduction is $100 \times (\mathrm{no\ NRTE} - \mathrm{NRTE})/\mathrm{no\ NRTE}$ computed per paired observation and then summarized by the median.}
\label{tab:raytrace_postproc_quality_250k}
\begin{tabular}{lrrrrr}
\toprule
\textbf{Metric} & \textbf{$N$} & \textbf{Median no NRTE} & \textbf{Median NRTE} & \textbf{Median reduction $\uparrow$ (\%)} & \textbf{Improved rows $\uparrow$ (\%)} \\
\midrule
SPO67 / Airy diameter & 30{,}415 & 168.24 & 22.95 & 81.1 & 89.7 \\
SPO100 / Airy diameter & 30{,}415 & 285.18 & 57.59 & 71.8 & 85.5 \\
EE50 / Airy diameter & 30{,}412 & 15.64 & 11.50 & 31.0 & 75.9 \\
EE90 / Airy diameter & 30{,}408 & 22.37 & 20.08 & 23.8 & 73.2 \\
Nominal performance & 31{,}776 & 11.76 & 2.73 & 70.9 & 81.6 \\
\bottomrule
\end{tabular}
\end{minipage}
\end{center}
